\chapter{Chimie supramoléculaire}\label{ch:b3:supramolecular}

En 1962, Charles Pedersen, qui cherchait autre chose, isola un polyéther
cyclique qui dissolvait le permanganate de potassium dans le benzène et le
teintait en violet. Le cycle s'était enroulé autour de l'ion potassium et
avait caché sa charge derrière une surface hydrocarbonée, entraînant le
permanganate comme contre-ion. La chimie au-delà de la molécule, celle des
assemblages maintenus par des forces plus faibles que les liaisons
covalentes, est née de cette solution violette. Ce chapitre mesure ces
forces, explique pourquoi certains \omterm{def:b3:supramolecular:supramolecular}{hôtes} choisissent si bien leurs \omterm{def:b3:supramolecular:supramolecular}{invités},
montre comment se mesurent les constantes d'association, et se termine par
des molécules enfilées et entrelacées en machines.

\begin{recall}
Le volume de première année a traité des forces intermoléculaires, de la
liaison hydrogène, de la solvatation, des molécules hydrophobes, des
complexes, des constantes de formation et des chélates~; le volume de
deuxième année, du potentiel chimique, des enthalpies libres et de
l'ajustement par moindres carrés. Le \cref{ch:b3:advanced-nmr} a décrit
l'échange rapide et lent et la coalescence, le \cref{ch:b3:rate-theories}
l'\omterm{thm:b3:rate-theories:eyring}{équation d'Eyring}, et le \cref{ch:b3:partition-functions} l'entropie de
translation d'un gaz.
\end{recall}

\section{Interactions non covalentes}

\begin{definition}[Chimie supramoléculaire]\label{def:b3:supramolecular:supramolecular}
La \emph{chimie supramoléculaire}\index{chimie supramoleculaire@chimie supramoléculaire} est la chimie
des assemblages de molécules maintenus par des interactions non covalentes.
Dans un \emph{complexe hôte--invité}\index{complexe hote--invite@complexe hôte--invité}, l'\emph{hôte}\index{hote@hôte}
est la plus grosse molécule, munie d'une cavité ou d'une fente et de sites
de liaison qui convergent vers elle, et l'\emph{invité}\index{invite@invité} la plus petite
molécule ou l'ion qu'elle fixe.
\end{definition}

Les interactions sont celles du volume de première année, ion--dipôle,
dipôle--dipôle, liaisons hydrogène et forces de London, plus quelques-unes
qui ont leur propre nom.

\begin{definition}[Interactions non covalentes]\label{def:b3:supramolecular:interactions}
L'\emph{empilement $\pi$}\index{empilement pi@empilement $\pi$} est l'attraction entre des cycles
aromatiques plans placés face à face, en général décalés ou inclinés. Une
\emph{interaction cation--$\pi$}\index{interaction cation--pi@interaction cation--$\pi$} est l'attraction entre un cation et
la face riche en électrons d'un cycle aromatique. Une \emph{liaison halogène}\index{liaison halogene@liaison halogène}
est l'attraction entre l'extrémité pauvre en électrons d'un atome d'halogène
lié de façon covalente, à l'opposé de sa liaison, et une base de Lewis.
L'\emph{effet hydrophobe}\index{effet hydrophobe} est la tendance des
surfaces apolaires à se rassembler dans l'eau, due surtout à la libération
des molécules d'eau ordonnées qui les entourent.
\end{definition}

Les interactions ion--dipôle et cation--$\pi$ sont les plus fortes, et
dépendent beaucoup du solvant qui entre en compétition pour l'ion~; viennent
ensuite les liaisons hydrogène, puis l'empilement $\pi$ et les \omterm{def:b3:supramolecular:interactions}{liaisons halogène}, puis les forces de
London, faibles une à une mais sommées sur de grandes surfaces de contact.
Dans l'eau, l'\omterm{def:b3:supramolecular:interactions}{effet hydrophobe} domine souvent~: un \omterm{def:b3:supramolecular:supramolecular}{hôte} conçu pour fixer un
\omterm{def:b3:supramolecular:supramolecular}{invité} par liaisons hydrogène dans le chloroforme peut ne pas le fixer du
tout dans l'eau, dont les propres liaisons hydrogène entrent en
compétition.

\section{Reconnaissance moléculaire}

\begin{definition}[Reconnaissance moléculaire]\label{def:b3:supramolecular:recognition}
La \emph{reconnaissance moléculaire}\index{reconnaissance moleculaire@reconnaissance moléculaire} est la fixation
sélective d'un \omterm{def:b3:supramolecular:supramolecular}{invité} par un \omterm{def:b3:supramolecular:supramolecular}{hôte}. Elle repose sur la \emph{complémentarité}\index{complementarite@complémentarité},
correspondance de taille, de forme et de sites de liaison entre l'\omterm{def:b3:supramolecular:supramolecular}{hôte} et
l'\omterm{def:b3:supramolecular:supramolecular}{invité}, et elle est renforcée par la \emph{préorganisation}\index{preorganisation@préorganisation}, propriété
d'un \omterm{def:b3:supramolecular:supramolecular}{hôte} dont les sites de liaison sont déjà maintenus dans la géométrie
du complexe avant la fixation.
\end{definition}

La constante d'association $K_a = [\mathrm{HG}]/([\mathrm H][\mathrm G])$ est une constante de formation au sens
du volume de première année, et $\Delta G^\circ = -RT\ln K_a$. Un \omterm{def:b3:supramolecular:supramolecular}{hôte} flexible doit figer sa
conformation lors de la fixation, ce qui coûte de l'entropie~; un \omterm{def:b3:supramolecular:supramolecular}{hôte}
rigide et préorganisé a payé ce prix lors de sa synthèse.

\begin{definition}[Effets chélate et macrocyclique]\label{def:b3:supramolecular:macrocyclic}
L'\emph{effet chélate}\index{effet chelate@effet chélate} est la plus grande stabilité d'un
complexe formé avec un ligand polydenté qu'avec le nombre équivalent de
ligands monodentés portant les mêmes atomes donneurs. L'\emph{effet macrocyclique}\index{effet macrocyclique}
est la stabilisation supplémentaire obtenue lorsque le ligand polydenté est
un cycle plutôt qu'une chaîne ouverte.
\end{definition}

\begin{proposition}[L'effet chélate est surtout entropique]\label{prop:b3:supramolecular:chelate-entropy}
Dans $\mathrm{[M(NH_3)_6]^{2+} + 3\,en \rightleftharpoons [M(en)_3]^{2+} + 6\,NH_3}$, où en désigne l'éthane-1,2-diamine, les six mêmes
liaisons métal--azote sont présentes des deux côtés, et la réaction est
menée par l'augmentation du nombre de molécules libres.
\end{proposition}

\begin{proof}[Justification]
Les liaisons formées et rompues sont de même nature, si bien que $\Delta_rH^\circ$ est faible. Le nombre de
molécules indépendantes passe de quatre à sept~: trois degrés de liberté de
translation de plus d'une molécule entière, valant chacun quelques dizaines
de joules par kelvin et par mole en solution (\cref{ch:b3:partition-functions}), donc $\Delta_rS^\circ > 0$ et $K > 1$.
Du point de vue cinétique, une fois une extrémité d'un chélate fixée,
l'autre est maintenue près du métal à une concentration effective élevée,
et se refixe aussitôt si elle se détache.
\end{proof}

\begin{proposition}[Compensation enthalpie--entropie]\label{prop:b3:supramolecular:compensation}
Dans une série de \omterm{def:b3:supramolecular:supramolecular}{complexes hôte--invité} apparentés, un $\Delta H^\circ$ de fixation plus négatif
s'accompagne en général d'un $\Delta S^\circ$ plus négatif, si bien que $\Delta G^\circ$ varie moins que chacun d'eux (c'est une
observation expérimentale).
\end{proposition}

Un complexe plus serré lie plus fortement mais restreint davantage les
mouvements de l'\omterm{def:b3:supramolecular:supramolecular}{hôte} et de l'\omterm{def:b3:supramolecular:supramolecular}{invité}~; une liaison hydrogène plus forte avec
l'\omterm{def:b3:supramolecular:supramolecular}{invité} signifie souvent une liaison plus faible avec l'eau qu'il a
déplacée. L'enthalpie libre, qui fixe la sélectivité, est la petite
différence de deux grands termes.

\begin{definition}[Hôtes]\label{def:b3:supramolecular:hosts}
Un \emph{éther-couronne}\index{ether-couronne@éther-couronne} est un polyéther macrocyclique, comme le
18-couronne-6, $\ce{(CH2CH2O)6}$. Un \emph{cryptand}\index{cryptand} est un polyéther
bicyclique à deux têtes de pont azotées, qui enferme un cation dans les
trois dimensions au sein d'un \emph{cryptate}\index{cryptate}. Une
\emph{cyclodextrine}\index{cyclodextrine} est un oligomère cyclique
d'unités glucose, en forme de cône tronqué, avec une cavité hydrophobe et
des groupes hydroxy sur ses bords. Un \emph{calixarène}\index{calixarene@calixarène} est un
macrocycle en forme de coupe formé d'unités phénol reliées par des ponts
méthylène.
\end{definition}

\begin{omfigure}
\begin{tikzpicture}[font=\small]
  % 18-crown-6 with K+
  \foreach \k in {0,...,17} {
    \pgfmathsetmacro\ang{90 + 20*\k}
    \pgfmathsetmacro\nxt{110 + 20*\k}
    \draw[black!70] (\ang:1.55) -- (\nxt:1.55);
  }
  \foreach \k in {0,3,...,15} {
    \pgfmathsetmacro\ang{90 + 20*\k}
    \node[fill=white, inner sep=0.5pt, text=omThm] (o\k) at (\ang:1.55) {O};
    \draw[omThm!60, densely dashed] (0,0) -- (\ang:1.32);
  }
  \node[circle, fill=omDef!25, draw=omDef, inner sep=1pt] at (0,0) {$\ce{K+}$};
  \node[font=\scriptsize] at (0,-2.1) {18-couronne-6 avec $\ce{K+}$};
  % [2.2.2]cryptand with K+
  \begin{scope}[xshift=5.2cm]
    \node (nt) at (0,1.7) {N};
    \node (nb) at (0,-1.7) {N};
    \foreach \x/\n in {-1.25/a, 0/b, 1.25/c} {
      \node[text=omThm, fill=white, inner sep=0.5pt] (oa\n) at (\x,0.55) {O};
      \node[text=omThm, fill=white, inner sep=0.5pt] (ob\n) at (\x,-0.55) {O};
      \draw[black!70] (nt) to[out=-90+40*\x, in=90] (oa\n) -- (ob\n) to[out=-90, in=90-40*\x] (nb);
    }
    \node[circle, fill=omDef!25, draw=omDef, inner sep=1pt] at (0.62,0) {$\ce{K+}$};
    \node[font=\scriptsize] at (0,-2.4) {[2.2.2]cryptand~: cryptate de $\ce{K+}$};
  \end{scope}
\end{tikzpicture}

{\small À gauche~: le 18-couronne-6, un cycle de douze atomes de carbone
(sommets) et six atomes d'oxygène, avec un ion potassium en son centre lié
par les six oxygènes. À droite~: le [2.2.2]\omterm{def:b3:supramolecular:hosts}{cryptand}, trois chaînes de deux
oxygènes d'éther entre deux têtes de pont azotées, schématique~: le cation
est maintenu dans une cage tridimensionnelle, par six oxygènes et les deux
azotes.}
\end{omfigure}

Les éthers-couronnes sélectionnent les cations par leur taille~: le 18-couronne-6 fixe $\ce{K+}$ plus fortement que
$\ce{Na+}$, plus petit, qui ne peut pas atteindre les six oxygènes à la fois, ou que $\ce{Cs+}$, plus gros, qui
se place au-dessus du cycle. Les \omterm{def:b3:supramolecular:hosts}{cryptands}, plus préorganisés et
enveloppant complètement l'ion, fixent plus fortement encore et
sélectionnent plus nettement. Caché dans l'\omterm{def:b3:supramolecular:supramolecular}{hôte}, le cation entraîne son
anion dans les solvants apolaires, où l'anion, non solvaté, devient très
réactif~: c'est le principe de la catalyse par transfert de phase.

\section{Mesurer l'association}

\begin{proposition}[Échange rapide]\label{prop:b3:supramolecular:fast-exchange}
Lorsqu'un \omterm{def:b3:supramolecular:supramolecular}{hôte} s'échange entre ses états libre et lié beaucoup plus vite
que la différence de leurs fréquences de résonance, son signal RMN apparaît
à la moyenne
$\delta_{\mathrm{obs}} = x_{\mathrm{free}}\delta_{\mathrm{free}} + x_{\mathrm{bound}}\delta_{\mathrm{bound}}$, où les $x$ sont les fractions d'\omterm{def:b3:supramolecular:supramolecular}{hôte} dans chaque état.
\end{proposition}

\begin{proof}
Chaque molécule d'\omterm{def:b3:supramolecular:supramolecular}{hôte} visite les deux états de nombreuses fois pendant la
mesure~; sa fréquence de précession, moyennée dans le temps, est la moyenne
pondérée des deux, avec pour poids les fractions de temps passées dans
chacun, qui sont à l'équilibre les fractions de molécules dans chaque état
(\cref{ch:b3:advanced-nmr}).
\end{proof}

\begin{theorem}[Isotherme 1:1 exacte]\label{thm:b3:supramolecular:isotherm}
Pour $\mathrm{H + G \rightleftharpoons HG}$, avec les concentrations totales $H_0$ et $G_0$ et la constante d'association $K$,
\[
  [\mathrm{HG}] = \frac{1}{2}\Bigl[\Bigl(H_0 + G_0 + \frac1K\Bigr) - \sqrt{\Bigl(H_0 + G_0 + \frac1K\Bigr)^2 - 4H_0G_0}\,\Bigr].
\]
\end{theorem}

\begin{proof}
Avec $c = [\mathrm{HG}]$, $K = c/((H_0 - c)(G_0 - c))$, c'est-à-dire $c^2 - (H_0 + G_0 + 1/K)c + H_0G_0 = 0$. Des deux racines,
seule la plus petite est inférieure à la fois à $H_0$ et à $G_0$ (leur produit vaut $H_0G_0$ et leur somme
dépasse $H_0 + G_0$)~: c'est celle qui porte le signe moins.
\end{proof}

\begin{omfigure}
\begin{tikzpicture}
\begin{axis}[name=a, width=0.48\linewidth, height=5.5cm, xmin=0, xmax=6, ymin=0, ymax=1.05,
  axis lines=left, xlabel={équivalents d'invité, $G_0/H_0$}, ylabel={fraction d'hôte liée},
  tick label style={font=\scriptsize}, label style={font=\small},
  ]
\addplot[omMeth, thick] table[x=equiv, y=K4] {figdata/out/bachelor-3-binding-titration.dat};
\addplot[omDef, thick] table[x=equiv, y=K3] {figdata/out/bachelor-3-binding-titration.dat};
\addplot[omThm, thick] table[x=equiv, y=K2] {figdata/out/bachelor-3-binding-titration.dat};
\node[font=\scriptsize, omMeth, anchor=north] at (axis cs:4,0.94) {$K = 10^4$};
\node[font=\scriptsize, omDef, anchor=north] at (axis cs:4,0.68) {$K = 10^3$};
\node[font=\scriptsize, omThm, anchor=north] at (axis cs:4,0.27) {$K = 10^2$};
\end{axis}
\begin{axis}[at={(a.east)}, anchor=west, xshift=1.4cm, width=0.48\linewidth, height=5.5cm,
  xmin=0, xmax=1, ymin=0, ymax=1.05, axis lines=left,
  xlabel={fraction molaire d'hôte}, ylabel={$[\text{complexe}]/c_{\text{total}}$},
  tick label style={font=\scriptsize}, label style={font=\small}]
\addplot[omDef, thick] table[x=x, y=HG] {figdata/out/bachelor-3-binding-titration-b.dat};
\addplot[omThm, thick] table[x=x, y=HG2] {figdata/out/bachelor-3-binding-titration-b.dat};
\draw[black!35, dashed] (axis cs:0.5,0) -- (axis cs:0.5,0.55) (axis cs:0.3333,0) -- (axis cs:0.3333,0.38);
\node[font=\scriptsize, omDef, anchor=south] at (axis cs:0.5,0.52) {HG};
\node[font=\scriptsize, omThm, anchor=south] at (axis cs:0.25,0.36) {$\mathrm{HG_2}$};
\end{axis}
\end{tikzpicture}

{\small À gauche~: fraction d'un \omterm{def:b3:supramolecular:supramolecular}{hôte} liée au cours d'un titrage à $H_0 = \qty{1.0}{mmol/L}$ pour trois
constantes d'association (en $\unit{L.mol^{-1}}$)~: plus $K$ est grand, plus la cassure à un équivalent est nette~;
lorsque $KH_0 \gg 1$, la courbe indique seulement que $K$ est grand. À droite~: \omterm{def:b3:supramolecular:job}{diagrammes de
Job}, la concentration du complexe à concentration totale constante d'\omterm{def:b3:supramolecular:supramolecular}{hôte}
et d'\omterm{def:b3:supramolecular:supramolecular}{invité}, maximale pour une fraction d'\omterm{def:b3:supramolecular:supramolecular}{hôte} de 1/2 pour un complexe 1:1
et de 1/3 pour un complexe 1:2 (constantes du modèle).}
\end{omfigure}

\begin{method}[Constante d'association par titrage RMN]\label{met:b3:supramolecular:nmr-titration}
\begin{enumerate}
  \item Maintenir fixe la concentration d'\omterm{def:b3:supramolecular:supramolecular}{hôte} $H_0$ et ajouter l'\omterm{def:b3:supramolecular:supramolecular}{invité},
        de 0 à plusieurs équivalents~; enregistrer un signal de l'\omterm{def:b3:supramolecular:supramolecular}{hôte} en
        chaque point (échange rapide).
  \item Modèle~: $\delta = \delta_{\mathrm{free}} + (\delta_{\mathrm{bound}} - \delta_{\mathrm{free}})[\mathrm{HG}]/H_0$, avec $[\mathrm{HG}]$ tiré de l'isotherme exacte.
  \item Ajuster $K$ et $\delta_{\mathrm{bound}}$ par moindres carrés non linéaires, en minimisant la somme des carrés des
        résidus~; tirer les écarts types de la courbure de cette somme (la
        matrice de covariance).
  \item Vérifier l'absence de tendance dans les résidus (mauvaise
        stœchiométrie) et que la fraction liée couvre à peu près 20 à
        80\,\% (les données contraignent alors $K$).
\end{enumerate}
\end{method}

\begin{definition}[Diagramme de Job]\label{def:b3:supramolecular:job}
Un \emph{diagramme de Job}\index{diagramme de Job} (méthode des variations
continues) est le tracé de la concentration d'un complexe, ou d'un signal
qui lui est proportionnel, en fonction de la fraction molaire de l'un des
constituants, à concentration totale constante des deux.
\end{definition}

\begin{proposition}[Maximum d'un diagramme de Job]\label{prop:b3:supramolecular:job}
Pour un complexe unique $\mathrm{H_nG_m}$, le \omterm{def:b3:supramolecular:job}{diagramme de Job} est maximal pour la fraction molaire d'\omterm{def:b3:supramolecular:supramolecular}{hôte}
$x = n/(n + m)$, quelle que soit la constante d'association.
\end{proposition}

\begin{proof}
Avec la concentration totale $T$, $H_t = xT$, $G_t = (1 - x)T$, le complexe $c$, les espèces libres $[\mathrm H] = H_t - nc$,
$[\mathrm G] = G_t - mc$, et $c = \beta[\mathrm H]^n[\mathrm G]^m$. En dérivant par rapport à $x$ au maximum, où
$\dd c/\dd x = 0$~: $0 = \beta\bigl(n[\mathrm H]^{n-1}[\mathrm G]^m T - m[\mathrm H]^n[\mathrm G]^{m-1}T\bigr)$, donc $n[\mathrm G] = m[\mathrm H]$. En substituant,
$n(G_t - mc) = m(H_t - nc)$, donc $nG_t = mH_t$, c'est-à-dire $n(1 - x) = mx$~: $x = n/(n + m)$.
\end{proof}

\begin{method}[Réaliser un diagramme de Job]\label{met:b3:supramolecular:job}
\begin{enumerate}
  \item Préparer des solutions mères équimolaires d'\omterm{def:b3:supramolecular:supramolecular}{hôte} et d'\omterm{def:b3:supramolecular:supramolecular}{invité}.
  \item Les mélanger dans des proportions $x$ allant de 0 à 1, à volume
        total constant.
  \item Mesurer une propriété proportionnelle au complexe (une absorbance
        corrigée des espèces libres, ou $x\,\Delta\delta$ en RMN).
  \item Lire la stœchiométrie sur la position du maximum.
\end{enumerate}
\end{method}

\begin{definition}[Titrage calorimétrique isotherme]\label{def:b3:supramolecular:itc}
Le \emph{titrage calorimétrique isotherme}\index{titrage calorimetrique isotherme@titrage calorimétrique isotherme}
(ITC) mesure la chaleur dégagée ou absorbée à chaque injection d'une
solution d'\omterm{def:b3:supramolecular:supramolecular}{invité} dans une solution d'\omterm{def:b3:supramolecular:supramolecular}{hôte} maintenue à température
constante.
\end{definition}

\begin{proposition}[Chaleur par injection]\label{prop:b3:supramolecular:itc}
En négligeant la dilution, la chaleur de la $i$-ième injection dans une cellule de volume $V_0$ vaut
$q_i = \Delta H^\circ V_0\bigl([\mathrm{HG}]_i - [\mathrm{HG}]_{i-1}\bigr)$.
\end{proposition}

\begin{proof}
La chaleur est l'enthalpie de réaction multipliée par la quantité de
complexe formée pendant l'injection, et cette quantité est la variation de $[\mathrm{HG}]$ multipliée par le volume de la cellule.
\end{proof}

\begin{method}[Lire une isotherme d'ITC]\label{met:b3:supramolecular:itc}
\begin{enumerate}
  \item Tracer la chaleur par mole d'\omterm{def:b3:supramolecular:supramolecular}{invité} injecté en fonction du rapport
        molaire.
  \item La hauteur du plateau avant saturation donne $\Delta H^\circ$~; le rapport molaire au
        point d'inflexion donne la stœchiométrie~; la raideur donne $K$.
  \item Ajuster les trois grandeurs avec l'isotherme exacte~; puis $\Delta G^\circ = -RT\ln K$ et
        $T\Delta S^\circ = \Delta H^\circ - \Delta G^\circ$.
\end{enumerate}
\end{method}

Une seule expérience d'ITC donne ainsi toute la signature thermodynamique
de l'association, enthalpie et entropie, qu'un titrage par RMN ou par
UV--visible ne donne qu'au moyen d'une série de températures.

\begin{inthelab}[Un titrage RMN]
Une solution de l'\omterm{def:b3:supramolecular:supramolecular}{hôte} dans un solvant deutéré, à environ une millimole par
litre, est placée dans un tube de RMN et son spectre enregistré. Une
solution concentrée de l'\omterm{def:b3:supramolecular:supramolecular}{invité}, préparée dans la solution d'\omterm{def:b3:supramolecular:supramolecular}{hôte} pour que
la concentration d'\omterm{def:b3:supramolecular:supramolecular}{hôte} reste constante, est ajoutée par petites portions à
la microseringue~; après chaque ajout, le tube est agité, équilibré à la
température de la sonde, et le spectre enregistré. On relève en chaque
point le déplacement d'un signal de l'\omterm{def:b3:supramolecular:supramolecular}{hôte} qui se déplace, puis on
l'ajuste.
\end{inthelab}

\section{Les hôtes}

Les \omterm{def:b3:supramolecular:hosts}{cyclodextrines}, produites par des \omterm{def:b3:complex-kinetics:enzyme}{enzymes} à partir de l'amidon,
comptent six, sept ou huit unités glucose ($\alpha$, $\beta$, $\gamma$). Leur cavité hydrophobe
accueille les parties apolaires des \omterm{def:b3:supramolecular:supramolecular}{invités} dans l'eau~: elles
solubilisent des médicaments, transportent des parfums et retirent le
cholestérol des membranes cellulaires. Les \omterm{def:b3:supramolecular:hosts}{calixarènes} et les
cucurbiturils en forme de citrouille, macrocycles rigides d'unités
glycolurile aux portails bordés de carbonyles, fixent très fortement les
cations et les cations organiques dans l'eau.

\begin{omfigure}
\begin{tikzpicture}[font=\small]
  \draw[black!70, fill=omDef!8] (-1.9,1.2) -- (-1.2,-1.2) arc[start angle=180, end angle=360, x radius=1.2, y radius=0.3]
    -- (1.9,1.2);
  \draw[black!70, fill=omDef!18] (0,1.2) ellipse[x radius=1.9, y radius=0.45];
  \draw[black!40, dashed] (-1.2,-1.2) arc[start angle=180, end angle=0, x radius=1.2, y radius=0.3];
  \node[draw=omThm, thick, regular polygon, regular polygon sides=6, minimum size=8mm, inner sep=0pt] at (0,0.2) {};
  \draw[omThm, thick] (0,0.62) -- (0,2.0);
  \node[font=\scriptsize, omThm, anchor=west] at (0.1,2.0) {invité};
  \node[font=\scriptsize, anchor=west] at (2.0,1.2) {grand bord~: OH secondaires};
  \node[font=\scriptsize, anchor=west] at (1.3,-1.25) {petit bord~: OH primaires};
  \node[font=\scriptsize, anchor=east] at (-1.75,0) {cavité hydrophobe};
\end{tikzpicture}

{\small Une \omterm{def:b3:supramolecular:hosts}{cyclodextrine}, dessinée comme le cône tronqué qu'elle forme,
avec un \omterm{def:b3:supramolecular:supramolecular}{invité} aromatique dans sa cavité. Les groupes hydroxy des deux
bords rendent l'extérieur soluble dans l'eau~; l'intérieur, tapissé de
groupes C--H et d'oxygènes glycosidiques, est apolaire.}
\end{omfigure}

\section{Auto-assemblage et machines moléculaires}

\begin{definition}[Auto-assemblage]\label{def:b3:supramolecular:self-assembly}
L'\emph{auto-assemblage}\index{auto-assemblage} est l'organisation
spontanée et réversible de composants en une structure définie, sous le
contrôle de l'information contenue dans les composants eux-mêmes~: leurs
formes et leurs sites de liaison.
\end{definition}

Des ions métalliques aux angles de coordination fixés et des ligands
pontants rigides s'assemblent en une seule étape en carrés, en cages et en
capsules, parce que chaque liaison peut s'ouvrir et se refermer jusqu'à
atteindre la structure la plus stable, sans tension ni site pendant.
L'appariement des bases dans l'ADN et le repliement des protéines relèvent
de la même idée, à plus grande échelle.

\begin{definition}[Molécules mécaniquement entrelacées]\label{def:b3:supramolecular:interlocked}
Une \emph{molécule mécaniquement entrelacée}\index{molecule mecaniquement entrelacee@molécule mécaniquement entrelacée}
possède des parties qui ne sont pas liées entre elles mais ne peuvent être
séparées sans rompre une liaison covalente. Un \emph{rotaxane}\index{rotaxane} est un anneau
enfilé sur un axe muni de bouchons encombrants à ses deux extrémités~; un
\emph{caténane}\index{catenane@caténane} est formé d'anneaux entrelacés.
\end{definition}

\begin{definition}[Synthèse par effet de gabarit]\label{def:b3:supramolecular:template}
Une \emph{synthèse par effet de gabarit}\index{synthese par effet de gabarit@synthèse par effet de gabarit} utilise une interaction
temporaire, souvent avec un ion métallique, pour maintenir des composants
réactifs dans la géométrie nécessaire à une réaction qui serait sinon
improbable, comme la fermeture d'un cycle autour d'un autre cycle.
\end{definition}

\begin{omfigure}
\begin{tikzpicture}[font=\small, >=Stealth]
  % step 1: two crescents around Cu+
  \draw[omDef, line width=2pt] (0.3,0) arc[start angle=0, end angle=300, radius=0.75];
  \draw[omThm, line width=2pt] (-0.3,0) arc[start angle=180, end angle=-120, radius=0.75];
  \fill[omProp] (0,0) circle (0.13);
  \node[font=\scriptsize] at (0,-1.35) {$\ce{Cu+}$ maintient deux};
  \node[font=\scriptsize] at (0,-1.7) {ligands ouverts croisés};
  \draw[->, thick] (1.4,0) -- (2.3,0) node[midway, above, font=\scriptsize] {\tiny fermeture};
  % step 2: closed interlocked rings with Cu
  \begin{scope}[xshift=3.6cm]
    \draw[omDef, line width=2pt] (-0.35,0) circle (0.75);
    \draw[omThm, line width=2pt] (0.35,0) circle (0.75);
    \draw[white, line width=5pt] ([shift={(52:0.75)}]-0.35,0) arc[start angle=52, end angle=72, radius=0.75];
    \draw[omDef, line width=2pt] ([shift={(48:0.75)}]-0.35,0) arc[start angle=48, end angle=76, radius=0.75];
    \fill[omProp] (0,0) circle (0.13);
    \node[font=\scriptsize] at (0,-1.35) {caténate};
  \end{scope}
  \draw[->, thick] (5.0,0) -- (5.9,0) node[midway, above, font=\scriptsize] {$-\ce{Cu+}$};
  \begin{scope}[xshift=7.2cm]
    \draw[omDef, line width=2pt] (-0.35,0) circle (0.75);
    \draw[omThm, line width=2pt] (0.35,0) circle (0.75);
    \draw[white, line width=5pt] ([shift={(52:0.75)}]-0.35,0) arc[start angle=52, end angle=72, radius=0.75];
    \draw[omDef, line width=2pt] ([shift={(48:0.75)}]-0.35,0) arc[start angle=48, end angle=76, radius=0.75];
    \node[font=\scriptsize] at (0,-1.35) {[2]caténane};
  \end{scope}
  % rotaxane
  \begin{scope}[xshift=10.2cm]
    \draw[black!70, line width=1.6pt] (-1.1,0) -- (1.1,0);
    \fill[black!60] (-1.25,0) circle (0.2) (1.25,0) circle (0.2);
    \draw[omThm, line width=2pt] (0,0) ellipse[x radius=0.18, y radius=0.55];
    \node[font=\scriptsize] at (0,-1.35) {[2]rotaxane};
  \end{scope}
\end{tikzpicture}

{\small Synthèse d'un \omterm{def:b3:supramolecular:interlocked}{caténane} par effet de \omterm{def:b3:supramolecular:template}{gabarit} (schéma)~: un ion
cuivre(I) maintient deux ligands dérivés de la phénanthroline croisés à
angle droit dans sa coordination tétraédrique~; refermer chacun en cycle
donne deux anneaux entrelacés autour du métal, et le retrait du cuivre
libère le \omterm{def:b3:supramolecular:interlocked}{caténane}. À droite~: un \omterm{def:b3:supramolecular:interlocked}{rotaxane}, anneau piégé sur un axe par
deux bouchons.}
\end{omfigure}

\begin{definition}[Machines moléculaires]\label{def:b3:supramolecular:machine}
Une \emph{machine moléculaire}\index{machine moleculaire@machine moléculaire} est un assemblage dont
les composants se déplacent les uns par rapport aux autres de façon
contrôlée en réponse à un stimulus (lumière, changement redox, changement
de pH), et que l'on peut faire travailler.
\end{definition}

Dans une navette moléculaire, l'anneau d'un \omterm{def:b3:supramolecular:interlocked}{rotaxane} se place sur l'une de
deux stations de l'axe~; un stimulus qui affaiblit sa fixation à la première
l'envoie vers la seconde, et le stimulus inverse le ramène. Les moteurs
rotatifs actionnés par la lumière, fondés sur des alcènes surencombrés,
tournent dans un seul sens, en alternant des isomérisations
photochimiques et des étapes thermiques qui ne peuvent pas s'inverser. La
vitesse du va-et-vient entre stations se mesure par RMN à température
variable, à partir de la coalescence des signaux des deux formes
(\cref{ch:b3:advanced-nmr}).

\begin{safety}{\ghs[0.62cm]{GHS07}}
% ledger: ghs4:18C6
18-Couronne-6~: nocif en cas d'ingestion, irritant pour la peau, les yeux
et les voies respiratoires. Les éthers-couronnes font passer des sels
métalliques dans les solvants organiques et à travers la peau~; ils sont
pesés et manipulés avec des gants.
\end{safety}

\begin{history}[La chimie supramoléculaire et ses machines]
Les éthers-couronnes de Charles Pedersen (1967), les \omterm{def:b3:supramolecular:hosts}{cryptands} de
Jean-Marie Lehn et les \omterm{def:b3:supramolecular:supramolecular}{hôtes} préorganisés de Donald Cram ont fondé le
domaine~; tous trois partagèrent le prix Nobel de chimie 1987. Jean-Pierre
Sauvage prépara en 1983 des \omterm{def:b3:supramolecular:interlocked}{caténanes} avec un bon rendement grâce au \omterm{def:b3:supramolecular:template}{gabarit} de
cuivre, Fraser Stoddart construisit des navettes à \omterm{def:b3:supramolecular:interlocked}{rotaxane}, et Ben Feringa
un moteur rotatif actionné par la lumière en 1999~; ils partagèrent le prix
Nobel de chimie 2016 pour la conception et la synthèse de \omterm{def:b3:supramolecular:machine}{machines
moléculaires}.
\end{history}

\section{Exercices}

\begin{exercise}[$\star$]\label{exo:b3:supramolecular:1}
Nommer l'interaction principale dans~: $\ce{K+}$ dans le 18-couronne-6~; deux cycles benzéniques face à face~;
$\ce{K+}$ au-dessus d'un cycle benzénique~; l'iodopentafluorobenzène avec la pyridine~; une paire
guanine--cytosine.
\end{exercise}

\begin{exercise}[$\star$]\label{exo:b3:supramolecular:2}
Calculer $\Delta G^\circ$ à \qty{298}{K} pour un \omterm{def:b3:supramolecular:supramolecular}{complexe hôte--invité} de constante $K_a = \qty{1.0e4}{L.mol^{-1}}$.
\end{exercise}

\begin{exercise}[$\star$]\label{exo:b3:supramolecular:3}
Dans le méthanol, $\log K = 6.1$ pour $\ce{K+}$ et 4.4 pour $\ce{Na+}$ avec le 18-couronne-6 (données de l'exercice).
Calculer la sélectivité et la différence des $\Delta G^\circ$.
\end{exercise}

\begin{exercise}[$\star$]\label{exo:b3:supramolecular:4}
Un \omterm{def:b3:supramolecular:job}{diagramme de Job} est maximal pour une fraction molaire d'\omterm{def:b3:supramolecular:supramolecular}{hôte} de 0.33. Quelle est la stœchiométrie~?
\end{exercise}

\begin{exercise}[$\star\star$]\label{exo:b3:supramolecular:5}
Un \omterm{def:b3:supramolecular:supramolecular}{hôte} à \qty{2.0}{mmol/L}, avec $K = \qty{1500}{L.mol^{-1}}$, $\delta_{\mathrm{free}} = \qty{3.560}{ppm}$ et $\delta_{\mathrm{bound}} = \qty{3.720}{ppm}$, reçoit un
équivalent d'\omterm{def:b3:supramolecular:supramolecular}{invité}. Calculer la fraction liée et le déplacement observé.
\end{exercise}

\begin{exercise}[$\star\star$]\label{exo:b3:supramolecular:6}
Pour le nickel(II), $\log\beta_6 = 8.6$ avec l'ammoniac et $\log\beta_3 = 18.3$ avec l'éthane-1,2-diamine (données de
l'exercice). Calculer la constante et $\Delta G^\circ$ à \qty{298}{K} de
$\ce{[Ni(NH3)6]^2+} + 3\,\mathrm{en} \rightleftharpoons \ce{[Ni(en)3]^2+} + \ce{6NH3}$, et dire ce qui la gouverne.
\end{exercise}

\begin{exercise}[$\star\star$]\label{exo:b3:supramolecular:7}
Une expérience d'ITC à \qty{298}{K} donne $K = \qty{2.0e5}{L.mol^{-1}}$ et $\Delta H^\circ = \qty{-40}{kJ/mol}$. Calculer $\Delta G^\circ$, $T\Delta S^\circ$ et
$\Delta S^\circ$.
\end{exercise}

\begin{exercise}[$\star\star$]\label{exo:b3:supramolecular:8}
Expliquer le rôle du cuivre(I), et de sa géométrie tétraédrique, dans la
synthèse de \omterm{def:b3:supramolecular:interlocked}{caténanes} de Sauvage. Comment le cuivre est-il retiré à la
fin~?
\end{exercise}

\begin{exercise}[$\star\star$]\label{exo:b3:supramolecular:9}
Les deux signaux d'une navette à \omterm{def:b3:supramolecular:interlocked}{rotaxane}, séparés de \qty{120}{Hz}, coalescent à \qty{300}{K} (données de
l'exercice). Calculer la constante de vitesse d'échange à la coalescence et
l'\omterm{def:b3:rate-theories:activation}{enthalpie libre d'activation} du va-et-vient.
\end{exercise}

\begin{exercise}[$\star\star\star$]\label{exo:b3:supramolecular:10}
Un médicament a une solubilité intrinsèque de \qty{0.10}{mmol/L} et forme un complexe 1:1 avec une
\omterm{def:b3:supramolecular:hosts}{cyclodextrine}, $K = \qty{500}{L.mol^{-1}}$ (données de l'exercice). Montrer que la solubilité totale dans une
solution de \omterm{def:b3:supramolecular:hosts}{cyclodextrine} totale $C_t$ vaut $S_0 + KS_0C_t/(1 + KS_0)$, et la calculer pour $C_t = \qty{10}{mmol/L}$.
\end{exercise}

\begin{exercise}[$\star\star\star$]\label{exo:b3:supramolecular:11}
Montrer à partir de l'isotherme exacte que, lorsque $KH_0 \gg 1$, la fraction liée est égale au
nombre d'équivalents jusqu'à un, puis reste égale à un. Pourquoi un tel
titrage ne donne-t-il qu'une borne inférieure de $K$~?
\end{exercise}

\begin{exercise}[$\star\star\star$]\label{exo:b3:supramolecular:12}
Pour l'\omterm{def:b3:supramolecular:supramolecular}{hôte} de l'\cref{exo:b3:supramolecular:5}, supposer l'association
contrôlée par la diffusion, $k_{\mathrm{on}} = \qty{1e9}{L.mol^{-1}.s^{-1}}$. Calculer $k_{\mathrm{off}}$ et montrer que l'échange est rapide à
l'échelle de temps de la RMN à \qty{400}{MHz}.
\end{exercise}

\section{Problème~: le benzène violet}

\begin{problem}[{Problème du week-end --- le benzène violet~: le potassium dans
un éther-couronne, l'extraction du permanganate dans un solvant organique,
une constante d'association tirée d'un titrage RMN, et pourquoi une
quantité catalytique de couronne suffit}]
\label{pb:b3:supramolecular:1}
% ledger: const:R, ghs:KMnO4, ghs:toluene, rion:K+VI, rion:Na+VI
Données du problème. Dans le méthanol à \qty{298}{K}, $\log K = 6.1$ pour $\ce{K+}$ et 4.4 pour $\ce{Na+}$ avec le
18-couronne-6. Extraction~: une phase aqueuse à \qty{0.10}{mol/L} de $\ce{KCl}$ et \qty{1.0}{mmol/L} de $\ce{KMnO4}$ est agitée
avec un volume égal de toluène contenant la couronne L~; la seule espèce
extraite est la paire d'ions $\ce{[KL]+}\ce{MnO4-}$, avec
$K_{\mathrm{ex}} = [\ce{[KL]+MnO4-}]_{\mathrm{org}}/([\ce{K+}]_{\mathrm{aq}}[\ce{MnO4-}]_{\mathrm{aq}}[\mathrm L]_{\mathrm{org}}) = \qty{1.0e5}{L^2.mol^{-2}}$, et la couronne reste dans le toluène. Titrage RMN d'un
\omterm{def:b3:supramolecular:hosts}{éther-couronne} (\qty{2.0}{mmol/L}) par un sel de sodium dans le méthanol deutéré, un signal $\ce{CH2}$
(ppm) en fonction du nombre d'équivalents de $\ce{Na+}$~:

\begin{center}\small
\begin{tabular}{l*{11}{c}}
\hline
équiv. & 0 & 0.25 & 0.5 & 0.75 & 1.0 & 1.5 & 2.0 & 3.0 & 4.0 & 6.0 & 8.0\\
$\delta$ & 3.560 & 3.590 & 3.612 & 3.635 & 3.649 & 3.674 & 3.688 & 3.697 & 3.704 & 3.708 & 3.714\\
\hline
\end{tabular}
\end{center}
% table: binding-titration-c

\medskip\noindent\textbf{Partie I --- Le complexe.}
\begin{enumerate}
  \item Dessiner le 18-couronne-6 et son complexe du potassium.
  \item Calculer le $\Delta G^\circ$ de fixation de $\ce{K+}$ et de $\ce{Na+}$.
  \item Calculer la sélectivité $\ce{K+}/\ce{Na+}$.
  \item Avec les rayons en coordinence six $r(\ce{K+}) = \qty{138}{pm}$ et $r(\ce{Na+}) = \qty{102}{pm}$, expliquer la sélectivité.
  \item Pourquoi le complexe est-il soluble dans le toluène alors que $\ce{KCl}$ ne l'est pas~?
  \item Pourquoi la couronne fixe-t-elle moins fortement dans l'eau que
        dans le méthanol~?
\end{enumerate}

\medskip\noindent\textbf{Partie II --- L'extraction.}
\begin{enumerate}[resume]
  \item Écrire l'équilibre d'extraction.
  \item En notant $y$ la concentration extraite, écrire l'équation vérifiée
        par $y$ lorsque
        $[\mathrm L]_0 = \qty{1.0}{mmol/L}$.
  \item La résoudre, et donner la fraction de permanganate extraite.
  \item Calculer la fraction pour $[\mathrm L]_0 = \qty{10}{mmol/L}$.
  \item La calculer pour $[\mathrm L]_0 = \qty{0.10}{mmol/L}$.
  \item Pourquoi le toluène est-il violet, et pourquoi le permanganate
        extrait est-il un oxydant plus fort que dans l'eau~?
\end{enumerate}

\medskip\noindent\textbf{Partie III --- Un titrage RMN.}
\begin{enumerate}[resume]
  \item Pourquoi un seul signal moyenné se déplace-t-il au lieu de deux
        signaux~?
  \item Écrire le modèle de $\delta$ en fonction de la concentration en \omterm{def:b3:supramolecular:supramolecular}{invité}.
  \item Ajuster $K$ et $\delta_{\mathrm{bound}}$ par moindres carrés~; donner $K$ avec son écart type.
  \item Donner la dispersion des résidus et commenter.
  \item Sur quel domaine d'équivalents la fraction liée est-elle comprise
        entre 20 et 80\,\%~?
  \item Pourquoi le même titrage ne permettrait-il pas de mesurer $K$ pour le potassium, $\log K = 6.1$~?
\end{enumerate}

\medskip\noindent\textbf{Partie IV --- Catalyse par transfert de phase.}
\begin{enumerate}[resume]
  \item Un alcène dans le toluène est oxydé par le permanganate extrait. Où
        la réaction a-t-elle lieu~?
  \item Que devient la couronne après la réaction, et pourquoi une
        quantité catalytique suffit-elle~?
  \item Pourquoi un excès de $\ce{KCl}$ dans l'eau est-il utile~?
  \item Quels catalyseurs moins chers font le même travail dans
        l'industrie~?
  \item Pourquoi utilise-t-on aujourd'hui le toluène là où Pedersen
        utilisait le benzène~?
  \item Énoncer le résultat~: la fraction de permanganate extraite pour
        $[\mathrm L]_0 = \qty{1.0}{mmol/L}$.
\end{enumerate}
\end{problem}
